\chapter{Kết luận}
\label{chap:conclusion}

\section{Kết quả đạt được}
Trong quá trình thực hiện khóa luận, nhóm phát triển đã hoàn thành xây dựng ứng dụng đọc truyện tranh trực tuyến MyBooksLibrary. Dưới đây là các kết quả chính đã đạt được:
\begin{itemize}
    \item \textbf{Sản phẩm đã:} 
    \begin{itemize}
        \item Xây dựng thành công ứng dụng với giao diện thân thiện, tương tác mượt mà.
        \item Hoàn thiện các tính năng cốt lõi: Khám phá, Tìm kiếm truyện, Đọc truyện (Reader) và Quản lý Thư viện cá nhân.
        \item Giải quyết tốt bài toán lưu trữ và đọc truyện ngoại tuyến (Offline) thông qua cơ chế đồng bộ và tải trước dữ liệu.
    \end{itemize}
    \item \textbf{Về quy trình phát triển:} Đã thiết lập thành công hệ thống Tích hợp liên tục (CI) kết hợp kiểm thử tự động, đảm bảo chất lượng và tính ổn định của mã nguồn trong suốt chu kỳ phát triển.
\end{itemize}

\section{Hạn chế của ứng dụng}
Mặc dù đã hoàn thiện các chức năng cốt lõi và đáp ứng được nhu cầu cơ bản của người dùng, ứng dụng vẫn còn một số điểm hạn chế về mặt tính năng và hệ thống:
\begin{itemize}
    \item \textbf{Phụ thuộc nguồn dữ liệu:} Hiện tại ứng dụng mới chỉ hỗ trợ một nguồn dữ liệu duy nhất là API của MangaDex. Nếu máy chủ này gặp sự cố hoặc thay đổi cấu trúc API, ứng dụng sẽ bị gián đoạn hoạt động.
    \item \textbf{Thiếu tính năng tương tác cộng đồng:} Dù đã có chức năng đánh giá (Review) qua Firebase, ứng dụng vẫn thiếu các tính năng giao tiếp trực tiếp như diễn đàn thảo luận, chia sẻ truyện hay bình luận chi tiết theo từng chương.
    \item \textbf{Chưa tối ưu hóa hoàn toàn dung lượng lưu trữ:} Tính năng tải truyện về đọc ngoại tuyến lưu trữ ảnh chất lượng nguyên bản chưa qua nén, dẫn đến việc chiếm dụng khá nhiều không gian bộ nhớ của thiết bị nếu người dùng tải số lượng lớn truyện.
\end{itemize}

\section{Hướng phát triển tương lai}
Để sản phẩm trở nên hoàn thiện và có thể triển khai thực tế tới đông đảo người dùng, nhóm phát triển đề xuất một số hướng đi trong tương lai:
\begin{itemize}
    \item \textbf{Xây dựng kiến trúc Đa nguồn (Multi-source):} Bổ sung các module theo cơ chế Plugin để cho phép người dùng chuyển đổi và lấy dữ liệu từ nhiều trang web/API truyện tranh khác nhau, giảm thiểu rủi ro phụ thuộc vào một máy chủ.
    \item \textbf{Mở rộng hệ sinh thái cộng đồng:} Nâng cấp hệ thống đám mây (Firebase) để hỗ trợ bình luận (Comment) chuyên sâu từng chương, tính năng chia sẻ và tương tác trực tiếp giữa các độc giả.
    \item \textbf{Hệ thống gợi ý thông minh (Recommendation System):} Ứng dụng các mô hình học máy (Machine Learning) để phân tích sở thích, thói quen đọc truyện của cá nhân, từ đó tự động đưa ra các đề xuất truyện phù hợp.
    \item \textbf{Tối ưu hóa quản lý dung lượng tự động:} Cho phép người dùng tùy chọn mức chất lượng nén ảnh khi tải về ngoại tuyến, đồng thời tích hợp tính năng tự động dọn dẹp các chương truyện đã đọc xong hoặc bộ nhớ đệm (cache) sau một khoảng thời gian quy định.
\end{itemize}